\documentclass[12pt]{article}
\usepackage[utf8]{inputenc}
\usepackage[T1]{fontenc}
\usepackage{hyperref}
\title{East Futuna Language Resources for K-12 Curriculum Development}
\author{Compiled by Claude Sonnet 5.5 (claude-sonnet-5-5) for Andr\'as Kornai}
\date{}
\begin{document}
\maketitle

\section{Introduction}
East Futuna (ISO 639-3: fud; Glottocode: east2447) is an Austronesian (Polynesian) language of Wallis and Futuna, with roughly 3,600 speakers. It is not known to be taught in schools. Searches found few resources, mostly reference records, one scripture translation and crowd-sourced online dictionaries. Related West Futunan (Vanuatu) is a different language, so check which variety a resource covers.

\section{Primary Educational Resources}
No primers, textbooks or school materials in East Futuna were found in the search results.

\section{Linguistic and Reference Materials}
\begin{itemize}
\item CRDO/Cocoon language record for East Futuna, which cites Claire Moyse-Faurie's 1997 grammar: \url{https://cocoon.huma-num.fr/exist/crdo/lexvo.xql?code=fud}
\item UN Digital Library record for a \emph{Dictionnaire futunien-fran\c{c}ais} (1993); publisher and access not confirmed: \url{https://digitallibrary.un.org/record/1284621?ln=fr}
\item Glosbe crowd-sourced East Futuna word pairs (accuracy varies): \url{https://glosbe.com/fud/lib}
\item Verbix East Futuna verb list (short): \url{https://docs.verbix.com/Languages/FutunaEast}
\item Wikipedia article on Futunan: \url{https://www.Wikipedia.com/wiki/Futunan}
\end{itemize}

\section{Religious and Community Resources}
\begin{itemize}
\item Scripture Earth page for East Futuna, listing the Bible translation \emph{Tosi Tapu}: \url{https://www.scriptureearth.org/00eng.php?iso=fud}
\end{itemize}

\section{Digital and Audio Resources}
Scripture Earth also lists Global Recordings Network audio for this language (see the link above). No other audio, video or corpus was found.

\section{Curriculum Development Context}
Little usable curriculum material was found. A grammar and a dictionary exist in print, but their availability was not confirmed. A curriculum would need new materials, built with the community and with the scholarly literature cited above.
\end{document}